% Primary and supplementary source basis: :contentReference[oaicite:5]{index=5} :contentReference[oaicite:6]{index=6}

\section{Complex Variables}

\begin{conceptbox}
    A \textbf{complex variable} is a variable of the form
    \[
        z=x+iy,
    \]
    where both $x$ and $y$ may vary independently over subsets of $\mathbb{R}$.
\end{conceptbox}

Unlike a real variable, which moves along a line, a complex variable moves in a plane. This makes the behavior of complex functions richer and also more restrictive when differentiability is required.

\subsection*{Motivation / Intuition}

In real calculus, a variable changes along one direction. In complex analysis, a variable can approach a point from infinitely many directions. This extra geometric freedom is the reason the theory becomes both more subtle and more powerful.

\subsection*{Formal Definitions}

\begin{conceptbox}
    A \textbf{neighborhood} of $z_0$ is a disk
    \[
        \{z: |z-z_0|<\varepsilon\}, \qquad \varepsilon>0.
    \]

    A \textbf{deleted neighborhood} of $z_0$ is
    \[
        \{z: 0<|z-z_0|<\varepsilon\}.
    \]

    A set is \textbf{open} if every point in it has a neighborhood contained in the set.

    A set is \textbf{connected} if any two points in it can be joined by a path lying entirely in the set.

    A \textbf{domain} is an open, connected set.
\end{conceptbox}

These ideas are needed because analyticity is not defined at just one isolated point. It is defined in neighborhoods and domains.

\subsection*{Geometric Interpretation}

\begin{itemize}
    \item The set
          \[
              |z|<1
          \]
          is the open unit disk.
    \item The set
          \[
              |z|\le 1
          \]
          contains the boundary circle and is not open.
    \item The set
          \[
              0<|z|<1
          \]
          is punctured at the origin but is still connected.
\end{itemize}

\subsection*{Key Ideas}

\begin{conceptbox}
    Complex analysis is local. To say a function behaves well at a point, we usually require good behavior not just at that point but in some neighborhood around it.
\end{conceptbox}

This is why concepts such as open sets and domains are not merely formal language. They are essential.

\begin{examplebox}
    \textbf{Example 1: Identify a domain}

    Determine whether
    \[
        D=\{z: |z|<2\}
    \]
    is a domain.

    The set consists of all points inside the circle of radius $2$ centered at the origin.

    It is open because each interior point has a small disk around it that still lies inside $D$.

    It is connected because any two points inside the disk can be joined by a line segment lying inside the disk.

    Hence,
    \[
        D
    \]
    is a domain.
\end{examplebox}

\begin{examplebox}
    \textbf{Example 2: A set that is not a domain}

    Consider
    \[
        S=\{z: |z|<1\}\cup\{z: |z-4|<1\}.
    \]

    This is the union of two disjoint open disks.

    It is open, because each disk is open.

    However, it is not connected: a point in the first disk cannot be joined to a point in the second disk by a path staying entirely inside $S$.

    Therefore, $S$ is not a domain.
\end{examplebox}

\begin{noteBox}
    In engineering, domains arise naturally when discussing where a transfer function is valid, where a logarithm is single-valued, or where a conformal map can be used safely.
\end{noteBox}

\subsection*{Common Mistakes / Notes}

\begin{itemize}
    \item A set can be open without being connected.
    \item A set can be connected even if it has a hole.
    \item Analyticity is always discussed with respect to a domain, not just isolated points.
\end{itemize}